% R(OC,sigma) crossing lambda^R twice - the multiplicity after Proposition 9.
% GENERATED BY generate_figures.py - edit that file, not this one.
% OC*_1=0.2130 (dR/dOC=+0.4612), OC*_2=0.4290 (dR/dOC=-0.5799)
% c=0.0600; OC runs over [0.0654, 0.4874], R over [0.1045, 0.2306]
\pgfplotsset{ocacc/.style={
  width=12.0cm, height=7.4cm,
  xmin=0.04, xmax=0.55, ymin=0.075, ymax=0.266,
  axis x line=bottom, axis y line=left,
  axis line style={-{Stealth[length=6pt]},black}, clip=false,
  xlabel={$OC$}, ylabel={$\mathcal{R}(OC,\sigma)$},
  every axis x label/.style={at={(current axis.south east)},anchor=west,xshift=9pt},
  every axis y label/.style={at={(current axis.north west)},anchor=south,yshift=7pt},
  tick label style={font=\footnotesize}, label style={font=\footnotesize},
  title style={font=\footnotesize, yshift=2pt},
  tick style={black},
  xtick={0.0600,0.2130,0.4290}, xticklabels={$c$,$OC^{*}_{1}$,$OC^{*}_{2}$},
  ytick={0.2000}, yticklabels={$\lambda^{R}$},
}}

\begin{figure}[htbp]
\centering
\begin{tikzpicture}
\begin{axis}[ocacc]
  \addplot[black,dashed,thick,samples=2,domain=0.04:0.55] {0.2000};
  \draw[gray,dashed] (axis cs:0.21303,0.075) -- (axis cs:0.21303,0.20000);
  \draw[gray,dashed] (axis cs:0.42901,0.075) -- (axis cs:0.42901,0.20000);
  \addplot[blue,very thick,smooth] coordinates {
(0.0654,0.1045) (0.0689,0.1074) (0.0726,0.1105) (0.0764,0.1135) (0.0803,0.1167)
(0.0843,0.1199) (0.0885,0.1232) (0.0928,0.1266) (0.0972,0.1300) (0.1018,0.1334)
(0.1066,0.1369) (0.1114,0.1404) (0.1164,0.1439) (0.1215,0.1475) (0.1268,0.1511)
(0.1322,0.1547) (0.1377,0.1583) (0.1434,0.1619) (0.1491,0.1655) (0.1550,0.1691)
(0.1610,0.1727) (0.1672,0.1762) (0.1734,0.1798) (0.1798,0.1833) (0.1862,0.1867)
(0.1928,0.1901) (0.1994,0.1935) (0.2061,0.1968) (0.2129,0.2000) (0.2198,0.2031)
(0.2268,0.2061) (0.2338,0.2090) (0.2409,0.2118) (0.2481,0.2144) (0.2553,0.2169)
(0.2625,0.2193) (0.2698,0.2214) (0.2771,0.2234) (0.2844,0.2252) (0.2918,0.2267)
(0.2991,0.2280) (0.3064,0.2291) (0.3138,0.2299) (0.3211,0.2304) (0.3284,0.2306)
(0.3357,0.2305) (0.3430,0.2301) (0.3502,0.2293) (0.3573,0.2283) (0.3644,0.2269)
(0.3715,0.2252) (0.3785,0.2232) (0.3854,0.2210) (0.3922,0.2184) (0.3990,0.2155)
(0.4056,0.2125) (0.4122,0.2092) (0.4187,0.2058) (0.4250,0.2023) (0.4313,0.1987)
(0.4375,0.1950) (0.4435,0.1913) (0.4494,0.1877) (0.4552,0.1842) (0.4609,0.1807)
(0.4665,0.1774) (0.4719,0.1742) (0.4772,0.1711) (0.4824,0.1682) (0.4874,0.1654)};
\end{axis}
\end{tikzpicture}
\caption{}
% \caption{The regular program acceptance rate $\mathcal{R}(OC,\sigma)$ against the opportunity cost $OC>c$ faced by the marginal top applicant, in units of $v^T$. It is not monotone, so the regular payline $\lambda^R$ is met at two opportunity costs $OC^*_1$ and $OC^*_2$, each of which induces an equilibrium in the sense of Definition \ref{def:equilibrium}. $F$ standard normal, $\sigma=0.5$; $G=\mathrm{Beta}(2,8)$ on $[-2.5,2.5]$; $\lambda^T=0.5$, $v^R/v^T=0.6$, $\lambda^R=0.2$, $c/v^R=0.1$.}
\label{fig:oc_acc}
\end{figure}
